\documentclass{article}

\usepackage[dblblindworkshop, final]{neurips_2026}

\usepackage[utf8]{inputenc}
\usepackage[T1]{fontenc}
\usepackage{hyperref}
\usepackage{url}
\usepackage{booktabs}
\usepackage{amsfonts}
\usepackage{amsmath}
\usepackage{microtype}
\usepackage{xcolor}
\usepackage{graphicx}
\usepackage{multirow}

\graphicspath{{figures_hpc_n750/}}

\workshoptitle{30562 --- Machine Learning and Artificial Intelligence}
\title{When Does Retrieval-Augmented Generation Help?}

\author{%
  Omar Jallouq \\
  Bocconi University \\
  \texttt{omar.jallouq@studbocconi.it} \\
  \And
  Filip Janczak \\
  Bocconi University \\
  \texttt{filip.janczak@studbocconi.it} \\
  \And
  Mohammad Aaqib Khatibi \\
  Bocconi University \\
  \texttt{mohammad.khatibi@studbocconi.it} \\
  \And
  Dhruv Singh \\
  Bocconi University \\
  \texttt{dhruv.singh@studbocconi.it}\\
}



\begin{document}

\maketitle

\begin{abstract}
Retrieval-Augmented Generation (RAG) augments language models with
non-parametric external memory, but its benefits depend on many
pipeline decisions that are rarely studied in controlled isolation.
We present a twelve-experiment ablation on TriviaQA ($n\!=\!750$
questions, pooled Wikipedia/web corpus) using Flan-T5-base as the
generator, covering retrieval method, chunk granularity, retrieval
depth, prompt template, oracle and cross-encoder re-ranking upper
bounds, external distractor robustness, controlled distraction at
fixed recall, forced retrieval miss, and generator scale.
Our main findings are:
(i) RAG provides a massive benefit over the parametric baseline
($\Delta\mathrm{EM}\!=\!{+0.551}$, Holm $p\!<\!10^{-4}$) because
Flan-T5-base genuinely lacks long-tail factual knowledge;
(ii) the generator, not the retriever, is the binding constraint at
this scale---retrieval method, embedding model, and cross-encoder
re-ranking are all statistically indistinguishable under Holm
correction;
(iii) with the answer guaranteed present, each additional hard topical
distractor monotonically reduces EM (0.673$\to$0.558 across
$N\!=\!0$--$8$ distractors); and
(iv) wrong-but-topical retrieved context primes parametric recall
($+0.102$ EM over an empty-context baseline), so a retrieval miss
is costly relative to correct retrieval but does not drag the model
below its parametric floor.
\end{abstract}

\section{Introduction}

Large language models (LLMs) encode substantial world knowledge during
pre-training, but this \emph{parametric memory} has well-known failure
modes: facts can be outdated, hallucinated, or simply absent for
long-tail entities [1].
Retrieval-Augmented Generation (RAG) addresses this by augmenting the
generator with a non-parametric corpus---retrieving relevant passages
at inference time and conditioning the model on them.

Despite RAG's strong empirical track record, its benefits are not
uniform and the contribution of individual components remains poorly
understood.
Retrieval can be noisy or fail entirely; irrelevant passages can
\emph{distract} the generator; and chunking strategy, embedding model,
and retrieval depth are typically tuned by intuition rather than
evidence.
Prior work establishes that RAG helps on knowledge-intensive tasks [1],
that how passages are integrated matters [2], and surveys the full
design space [3]---but causal component-level ablations under a
consistent evaluation protocol on the same benchmark are rare.

We address this gap with a controlled twelve-experiment study.
Our key contributions are:
\begin{itemize}
  \item A comprehensive ablation spanning retrieval family (BM25 /
    TF-IDF / Dense), embedding model, chunk granularity, retrieval
    depth~$k$, prompt template, oracle upper bound, cross-encoder
    re-ranking, external distractor robustness, and generator scale.
  \item Statistical evidence (Holm-corrected paired tests) that the
    \emph{generator} is the binding constraint: retrieval method,
    embedding model, and re-ranking are all indistinguishable within
    $\pm0.02$ EM---a \emph{positive} ``they tie'' conclusion, not a
    power-limited one.
  \item The first clean causal estimate of when RAG hurts: pinning
    recall at~1 and adding hard topical distractors causes a monotone
    11.5~pp EM drop (0.673$\to$0.558), entirely attributable to
    generator distraction.
  \item Evidence that topical retrieval misses prime parametric recall
    ($+0.102$ EM over an empty-context baseline, Holm $p\!<\!10^{-4}$),
    and that the RAG--No-RAG gap begins to shrink as model scale grows.
\end{itemize}


\section{Method}

\paragraph{System.}
We implement a standard retrieve-then-read pipeline.
A retriever selects the top-$k$ passages from a corpus; these are
concatenated into a prompt (question repeated before and after the
context for robustness to truncation); and the generator produces a
short answer.
We evaluate three retrieval families:
\textbf{BM25} (Okapi BM25 on a shared NFKD-normalised token stream),
\textbf{TF-IDF} (\texttt{TfidfVectorizer}, sublinear TF, same
tokeniser as BM25),
and \textbf{Dense} (sentence-transformers
\texttt{all-MiniLM-L6-v2}, 384-dim, cosine similarity).
For the embedding ablation (Exp.~10) and cross-encoder re-ranking
(Exp.~7) we additionally use \texttt{bge-small-en-v1.5},
\texttt{e5-small-v2}, and \texttt{ms-marco-MiniLM-L-6-v2}.
The generator is \textbf{Flan-T5-base} [5] ($\approx$250M parameters),
with greedy decoding, a 1024-token input budget, and middle-truncation
to preserve the answer cue.

\paragraph{Dataset and corpus.}
We use \textbf{TriviaQA} \texttt{rc}~[4], sampling 750 questions from
the validation split (seed 42).
TriviaQA was chosen because the primary references for this work---Lewis
et al.~[1] and Izacard and Grave~[2]---both evaluate on TriviaQA,
enabling direct comparison with those baselines.
TriviaQA provides Wikipedia entity pages and web search hits for each
question; pooling all documents across questions gives a realistic
retrieval setting where the gold evidence is always present but must
be located among $\sim$2{,}500 documents (including 2{,}000
topic-agnostic Simple English Wikipedia distractors added to widen the
noise floor).
Documents are chunked with a fixed word-window (default 128 words,
32-word overlap).
Gold answers are SQuAD-normalised (lowercase, strip articles/punctuation)
and a prediction is correct if it matches any gold alias.

\paragraph{Evaluation.}
We report \textbf{Exact Match (EM)}, \textbf{Token F1}, and
\textbf{Retrieval Recall@$k$} (a whole-word token-span match, avoiding
spurious partial-word substring hits).
All metrics carry 95\% percentile bootstrap CIs (1{,}000 resamples).
For pairwise comparisons we use \textbf{McNemar's exact test} on EM and
a \textbf{paired bootstrap} on score deltas, with
\textbf{Holm--Bonferroni correction} across the family of nine contrasts.
Paired tests cancel question-to-question difficulty variance, providing
substantially more power than comparing marginal CIs.

\section{Experiments and Results}

\paragraph{RAG vs.\ No-RAG (Exp.\ 5).}
RAG provides a massive benefit over the parametric floor
(Table~\ref{tab:core}, top):
$\Delta\mathrm{EM}\!=\!{+0.551}$, Holm $p\!<\!10^{-4}$ (discordant
419/6 questions).
The No-RAG EM of 0.049 confirms that Flan-T5-base genuinely lacks these
long-tail facts---the model produces confidently-shaped but incorrect
entity strings (e.g., \texttt{henry viii} for ``PM after Balfour'').
Of 750 questions, \textbf{411 benefit}, \textbf{319 tie}, and only
\textbf{6 are hurt} by retrieval; the 6 ``hurts'' are questions the
model already answered correctly from memory and where retrieved context
distracted it.

\begin{table}[t]
\centering
\caption{Core results. All experiments use Flan-T5-base, chunk=128
  words, $k\!=\!5$, instructed prompt, $n\!=\!750$ (Exp.~10:
  $n\!=\!400$). ``Holm $p$'' is the McNemar $p$-value adjusted across
  the nine-contrast family; ``(ref)'' marks the comparison baseline.}
\label{tab:core}
\small
\begin{tabular}{lcccc}
\toprule
Condition & EM & F1 & R@5 & Holm $p$ \\
\midrule
\multicolumn{5}{l}{\textit{RAG vs.\ No-RAG (Exp.\ 5)}} \\
No-RAG (parametric only) & 0.049 & 0.100 & --- & \\
RAG (Dense, $k\!=\!5$) & \textbf{0.600} & \textbf{0.668} & 0.865 &
  $<\!10^{-4}$ \\
\midrule
\multicolumn{5}{l}{\textit{Retriever comparison (Exp.\ 1)}} \\
BM25   & 0.597 & 0.670 & 0.859 & 1.00 \\
TF-IDF & 0.603 & 0.673 & 0.852 & 1.00 \\
Dense (MiniLM) & 0.600 & 0.668 & 0.865 & (ref) \\
\midrule
\multicolumn{5}{l}{\textit{Embedding model ablation (Exp.\ 10)}} \\
Dense (MiniLM)  & \textbf{0.623} & \textbf{0.680} & 0.877 & (ref) \\
Dense (bge-small) & 0.620 & 0.673 & 0.905 & --- \\
Dense (e5-small)  & 0.610 & 0.663 & \textbf{0.915} & --- \\
\midrule
\multicolumn{5}{l}{\textit{Oracle upper bound (Exp.\ 6)}} \\
RAG (Dense, $k\!=\!5$)       & 0.600 & 0.668 & 0.865 & (ref) \\
Oracle (gold forced in) & 0.619 & 0.684 & 0.999 & 0.0007 \\
\bottomrule
\end{tabular}
\end{table}

\paragraph{Retrieval pipeline: method, chunk size, and depth (Exps.\ 1--4).}
All three retrieval families are statistically indistinguishable
(Table~\ref{tab:core}, middle; Holm $p\!=\!1.0$ for all three pairs):
Dense retrieval achieves the highest Recall@5 (0.865) but not the
highest EM, confirming that residual error is generator-side.
At a fixed 1024-token budget (Exp.~2), smaller chunks win: 64-word
chunks ($\approx$8 packed per prompt) achieve EM~0.596 vs.\ 0.517 for
512-word chunks ($\approx$1.2 per prompt), since large chunks risk
losing the entire context window on a single retrieval miss.
Retrieval depth~$k$ (Exp.~3, chunk\,=\,48 words) shows monotone
improvement through $k\!=\!10$ (EM: 0.433$\to$0.583) with no distraction
penalty---consistent with distraction only materialising once the answer
is already guaranteed present.
Prompt template (Exp.~4) has no significant effect: concise vs.\
instructed differ by only 0.016 EM (Holm $p\!=\!0.18$).

\paragraph{Retrieval quality ceiling: oracle, re-ranking, and embedding model (Exps.\ 6, 7, 10).}
Three experiments probe whether better retrieval quality can raise EM,
and all point to the same conclusion: the generator is the binding
constraint.
The oracle upper bound (Exp.~6, Table~\ref{tab:core}, bottom) forces the
gold chunk in for the 13\% of questions that Dense misses, yielding
only $+0.019$ EM (Holm $p\!=\!0.0007$)---small but reliable.
Cross-encoder re-ranking (Exp.~7, \texttt{ms-marco-MiniLM-L-6-v2},
top-50$\to$top-5) raises Recall@5 by 0.017 but the EM gain ($+0.017$,
Holm $p\!=\!0.73$) is not significant.
The embedding model ablation (Exp.~10, $n\!=\!400$) reveals the same
pattern: e5-small reaches Recall@5$\!=\!$0.915 (highest) yet
EM$\!=\!$0.610 (lowest), while MiniLM has Recall@5$\!=\!$0.877 but
EM$\!=\!$0.623 (highest)---a stronger retriever finds more gold chunks;
the generator extracts no better answers from them.

\paragraph{Distractor robustness and retrieval failure (Exps.\ 8, 9, 11).}
Experiment~8 sweeps the external distractor count
($\{0, 500, 2000, 5000\}$ Simple-Wikipedia articles): Recall@5 and EM
are exactly flat across a $2.8\times$ corpus-size range, confirming
that topic-agnostic noise never out-ranks gold pages.
To isolate generator distraction from retrieval failure, Experiment~9
pins the gold chunk at rank~1 (Recall@$k\!=\!1$ throughout) and adds
$N\in\{0,1,2,4,8\}$ hard topical non-gold chunks.
EM drops \textbf{monotonically} from 0.673 ($N\!=\!0$, single-passage
ceiling) to 0.558 ($N\!=\!8$)---a fall of 11.5~pp attributable entirely
to distraction (Figure~\ref{fig:main}, left).
Experiment~11 (forced retrieval miss, Recall@$k\!=\!0$ by construction)
closes the analysis from the other side: wrong-but-topical context
scores EM~0.159---far below RAG (0.600) but well above no-context
baselines (0.049--0.057).
Decomposing: the instructed-prompt scaffold alone adds $+0.008$ EM over
the bare parametric baseline; topical-but-wrong context adds a further
$+0.102$ EM (Holm $p\!<\!10^{-4}$).
A retrieval miss therefore primes parametric recall rather than dragging
the model below its floor.

\begin{figure}[t]
  \centering
  \includegraphics[width=0.48\linewidth]{fig11_distraction.png}%
  \hfill%
  \includegraphics[width=0.48\linewidth]{fig15_generator_scale.png}
  \caption{%
    \textbf{Left (Exp.~9):} EM and F1 decline monotonically as hard
    topical distractors are added alongside a guaranteed gold chunk
    (Recall@$k\!=\!1$ throughout).
    The $N\!=\!0$ point (gold chunk only) is the single-passage
    generator ceiling.
    \textbf{Right (Exp.~12):} RAG vs.\ No-RAG EM across Flan-T5 scales.
    The No-RAG parametric floor rises monotonically; the RAG--No-RAG gap
    peaks at base and begins to shrink at large, confirming that
    retrieval helps most when parametric memory is weakest.}
  \label{fig:main}
\end{figure}

\paragraph{Generator scale (Exp.\ 12).}
Figure~\ref{fig:main} (right) sweeps Flan-T5 small/base/large (80M /
250M / 780M) on 150 questions.
The No-RAG parametric floor rises monotonically with scale
(0.027$\to$0.053$\to$0.113), while RAG rises more slowly
(0.467$\to$0.600$\to$0.633) toward the retrieval/generation ceiling.
The RAG--No-RAG gap \textbf{peaks at base} ($\Delta\!=\!0.547$) and
\textbf{begins to shrink at large} ($\Delta\!=\!0.520$), consistent with
retrieval helping most when parametric memory is weakest.
The large model is also the only one where RAG occasionally hurts
(3/150 questions) as growing parametric knowledge conflicts with retrieved
context.

\section{Conclusion}

We conducted a twelve-experiment controlled ablation of a RAG pipeline
on TriviaQA.
Four contrasts survive Holm correction:
RAG vs.\ No-RAG ($+0.551$ EM, the dominant effect);
both comparisons of forced-miss against no-context baselines
($+0.102$--$0.110$ EM, topical priming);
and Oracle vs.\ RAG ($+0.019$ EM, a small but reliable retrieval-failure
component first resolved at $n\!=\!750$).
Retrieval method, embedding model, and re-ranking are all
Holm-non-significant within $\pm0.02$ EM with large discordant
counts---a positive ``they tie'' result.
Experiment~9 provides the clearest causal picture of when RAG hurts:
each additional topical distractor reduces EM monotonically even when
the answer is guaranteed present.
Taken together, the results suggest that for small instruction-tuned
generators, RAG effort is better invested in the generator than in
retrieval quality.

\paragraph{Limitations.}
The study uses small encoder-decoders (up to 780M parameters) on a
corpus much smaller and less adversarial than open-domain Wikipedia.
A DPR model fine-tuned on TriviaQA would likely widen the sparse--dense
gap currently observed as a tie.
Much larger generators (LLaMA-class) would raise the parametric floor
further and shrink the RAG benefit; multi-hop questions and
adversarially constructed distractors are left to future work.

\section*{References}

\small

[1] Lewis, P., Perez, E., Piktus, A., Petroni, F., Karpukhin, V.,
Goyal, N., \ldots Kiela, D.\ (2020). Retrieval-Augmented Generation for
Knowledge-Intensive NLP Tasks. \textit{NeurIPS 2020}.

[2] Izacard, G., \& Grave, E.\ (2021). Leveraging Passage Retrieval with
Generative Models for Open Domain Question Answering.
\textit{EACL 2021}.

[3] Gao, Y., Xiong, Y., Gao, X., Jia, K., Pan, J., Bi, Y., \ldots Wang, H.\
(2024). Retrieval-Augmented Generation for Large Language Models:
A Survey. \textit{arXiv:2312.10997}.

[4] Joshi, M., Choi, E., Weld, D., \& Zettlemoyer, L.\ (2017). TriviaQA:
A Large Scale Distantly Supervised Challenge Dataset for Reading
Comprehension. \textit{ACL 2017}.

[5] Chung, H.~W., Hou, L., Longpre, S., Zoph, B., Tay, Y., Fedus, W.,
\ldots Wei, J.\ (2022). Scaling Instruction-Finetuned Language Models.
\textit{arXiv:2210.11416}.

[6] Robertson, S., \& Zaragoza, H.\ (2009). The Probabilistic Relevance
Framework: BM25 and Beyond.
\textit{Foundations and Trends in Information Retrieval}, 3(4).

\appendix

\section{Additional Figures}

Figures~\ref{fig:app_retriever}--\ref{fig:app_scale} show all
fifteen experiment plots from the report-scale run ($n\!=\!750$,
$n\!=\!400$ for Exp.~10, $n\!=\!150$ for Exp.~12) on GPU (HPC).
Error bars are 95\% percentile bootstrap CIs (1{,}000 resamples).

\begin{figure}[h]
  \begin{minipage}{0.48\textwidth}
    \centering
    \includegraphics[width=\linewidth]{fig1_retriever_comparison.png}
    \caption{Exp.~1 --- Retriever comparison (BM25 / TF-IDF / Dense).
      All three are statistically indistinguishable (Holm $p\!=\!1.0$
      for all pairs).}
    \label{fig:app_retriever}
  \end{minipage}
  \hfill
  \begin{minipage}{0.48\textwidth}
    \centering
    \includegraphics[width=\linewidth]{fig2_chunk_size.png}
    \caption{Exp.~2 --- Chunk-size ablation at a fixed 1024-token budget.
      Smaller chunks pack more passages and achieve higher recall and EM.}
    \label{fig:app_chunk}
  \end{minipage}
\end{figure}

\begin{figure}[h]
  \begin{minipage}{0.48\textwidth}
    \centering
    \includegraphics[width=\linewidth]{fig3_k_values.png}
    \caption{Exp.~3 --- Retrieval depth sweep ($k\!\in\!\{1,3,5,10\}$,
      chunk\,=\,48 words).
      EM and Recall@$k$ both increase monotonically; no distraction penalty
      appears at standard retrieval.}
    \label{fig:app_k}
  \end{minipage}
  \hfill
  \begin{minipage}{0.48\textwidth}
    \centering
    \includegraphics[width=\linewidth]{fig4_prompt_template.png}
    \caption{Exp.~4 --- Prompt template: concise vs.\ instructed.
      Difference is 0.016 EM (Holm $p\!=\!0.18$, non-significant).}
    \label{fig:app_prompt}
  \end{minipage}
\end{figure}

\begin{figure}[h]
  \begin{minipage}{0.48\textwidth}
    \centering
    \includegraphics[width=\linewidth]{fig5_rag_vs_no_rag.png}
    \caption{Exp.~5 --- RAG vs.\ No-RAG.
      $\Delta$EM\,=\,+0.551 (Holm $p\!<\!10^{-4}$; 411 helps / 6 hurts /
      319 ties out of 750 questions).}
    \label{fig:app_rag_vs_norag}
  \end{minipage}
  \hfill
  \begin{minipage}{0.48\textwidth}
    \centering
    \includegraphics[width=\linewidth]{fig6_error_analysis.png}
    \caption{Exp.~5 error analysis: per-question RAG helps / hurts / ties
      breakdown and scatter of Recall@5 vs.\ $\Delta$EM.}
    \label{fig:app_error}
  \end{minipage}
\end{figure}

\begin{figure}[h]
  \begin{minipage}{0.48\textwidth}
    \centering
    \includegraphics[width=\linewidth]{fig7_qualitative.png}
    \caption{Exp.~5 qualitative examples where RAG helps and where it hurts.}
    \label{fig:app_qual}
  \end{minipage}
  \hfill
  \begin{minipage}{0.48\textwidth}
    \centering
    \includegraphics[width=\linewidth]{fig8_oracle.png}
    \caption{Exp.~6 --- Oracle vs.\ RAG vs.\ No-RAG.
      Oracle gain: $+0.019$ EM (Holm $p\!=\!0.0007$), quantifying the
      small reliable retrieval-failure component.}
    \label{fig:app_oracle}
  \end{minipage}
\end{figure}

\begin{figure}[h]
  \begin{minipage}{0.48\textwidth}
    \centering
    \includegraphics[width=\linewidth]{fig9_rerank.png}
    \caption{Exp.~7 --- Cross-encoder re-ranking (e5 first stage,
      top-50$\to$top-5).
      Recall@5 rises by 0.017 but EM gain ($+0.017$) is not
      significant (Holm $p\!=\!0.73$).}
    \label{fig:app_rerank}
  \end{minipage}
  \hfill
  \begin{minipage}{0.48\textwidth}
    \centering
    \includegraphics[width=\linewidth]{fig10_distractor_sweep.png}
    \caption{Exp.~8 --- External distractor count sweep
      ($\{0, 500, 2000, 5000\}$ Simple-Wikipedia articles).
      Recall@5 and EM are exactly flat across a $2.8\times$
      corpus-size range: topic-agnostic distractors never
      out-rank gold pages.}
    \label{fig:app_distractor_sweep}
  \end{minipage}
\end{figure}

\begin{figure}[h]
  \begin{minipage}{0.48\textwidth}
    \centering
    \includegraphics[width=\linewidth]{fig11_distraction.png}
    \caption{Exp.~9 --- Controlled distraction (Recall@$k\!=\!1$
      throughout). Adding hard topical distractors drives EM monotonically
      from 0.673 ($N\!=\!0$) to 0.558 ($N\!=\!8$).}
    \label{fig:app_distraction}
  \end{minipage}
  \hfill
  \begin{minipage}{0.48\textwidth}
    \centering
    \includegraphics[width=\linewidth]{fig12_significance.png}
    \caption{Paired significance table (McNemar on EM, paired bootstrap
      on deltas; raw and Holm-adjusted $p$-values for all nine contrasts).}
    \label{fig:app_sig}
  \end{minipage}
\end{figure}

\begin{figure}[h]
  \begin{minipage}{0.48\textwidth}
    \centering
    \includegraphics[width=\linewidth]{fig13_embedding_models.png}
    \caption{Exp.~10 --- Embedding model ablation ($n\!=\!400$).
      e5-small achieves the highest Recall@5 (0.915) yet the
      lowest EM (0.610); MiniLM has the lowest recall but the
      highest EM (0.623).}
    \label{fig:app_embed}
  \end{minipage}
  \hfill
  \begin{minipage}{0.48\textwidth}
    \centering
    \includegraphics[width=\linewidth]{fig14_forced_miss.png}
    \caption{Exp.~11 --- Forced retrieval miss (Recall@$k\!=\!0$).
      Topical-but-wrong context adds $+0.102$ EM over the empty-context
      baseline (Holm $p\!<\!10^{-4}$) via parametric priming.}
    \label{fig:app_miss}
  \end{minipage}
\end{figure}

\begin{figure}[h]
  \centering
  \includegraphics[width=0.6\linewidth]{fig15_generator_scale.png}
  \caption{Exp.~12 --- Generator-scale sweep ($n\!=\!150$).
    No-RAG floor rises with scale (0.027$\to$0.053$\to$0.113);
    RAG--No-RAG gap peaks at base ($\Delta\!=\!0.547$) and shrinks
    at large ($\Delta\!=\!0.520$).}
  \label{fig:app_scale}
\end{figure}

\section{Complete Results Tables}

\begin{table}[h]
\centering
\caption{Exp.~2 --- Chunk size at fixed 1024-token budget (Dense,
  $n\!=\!750$, 30 candidates packed greedily).}
\small
\begin{tabular}{rcccc}
\toprule
Chunk (words) & Mean packed $k$ & EM & Token F1 & Recall@$k$ \\
\midrule
64  & 8.27 & 0.596 & 0.669 & 0.859 \\
128 & 4.29 & 0.584 & 0.653 & 0.855 \\
256 & 2.14 & 0.537 & 0.614 & 0.820 \\
512 & 1.19 & 0.517 & 0.584 & 0.759 \\
\bottomrule
\end{tabular}
\end{table}

\begin{table}[h]
\centering
\caption{Exp.~3 --- Retrieval depth $k$ (Dense, chunk\,=\,48 words,
  $n\!=\!750$; full range fits 1024-token budget untruncated).}
\small
\begin{tabular}{rccc}
\toprule
$k$ & EM & Token F1 & Recall@$k$ \\
\midrule
1  & 0.433 & 0.496 & 0.477 \\
3  & 0.525 & 0.591 & 0.683 \\
5  & 0.547 & 0.615 & 0.768 \\
10 & 0.583 & 0.649 & 0.855 \\
\bottomrule
\end{tabular}
\end{table}

\begin{table}[h]
\centering
\caption{Exp.~9 --- Controlled distraction (Recall@$k\!=\!1$ by
  construction; gold chunk pinned at rank~1, $n\!\approx\!750$,
  chunk\,=\,48 words).}
\small
\begin{tabular}{rccc}
\toprule
$N$ distractors & EM & Token F1 & Recall@$k$ \\
\midrule
0 & 0.673 & 0.745 & 1.00 \\
1 & 0.617 & 0.680 & 1.00 \\
2 & 0.590 & 0.656 & 1.00 \\
4 & 0.575 & 0.645 & 1.00 \\
8 & 0.558 & 0.629 & 1.00 \\
\bottomrule
\end{tabular}
\end{table}

\begin{table}[h]
\centering
\caption{Exp.~11 --- Forced retrieval miss ($n\!\approx\!750$,
  chunk\,=\,48 words; Recall@$k\!=\!0$ by construction for the
  forced-miss arm).}
\small
\begin{tabular}{lccc}
\toprule
Condition & EM & Token F1 & Recall@$k$ \\
\midrule
No-RAG (bare \texttt{Q:/A:} cue) & 0.049 & 0.100 & --- \\
Matched (instructed prompt, \emph{empty} context) & 0.057 & 0.101 & --- \\
Forced miss (instructed prompt, non-gold chunks) & 0.159 & 0.246 & 0.00 \\
RAG (Dense, $k\!=\!5$, Exp.~5) & 0.600 & 0.668 & 0.865 \\
\bottomrule
\end{tabular}
\end{table}

\begin{table}[h]
\centering
\caption{Exp.~12 --- Generator-scale sweep ($n\!=\!150$; trend
  demonstration on a representative subset).}
\small
\begin{tabular}{lcccc}
\toprule
Generator & No-RAG EM & RAG EM & $\Delta$EM &
  Helps / Hurts / Ties \\
\midrule
\texttt{flan-t5-small} ($\approx$80M)  & 0.027 & 0.467 & $+0.440$ &
  67 / 1 / 82 \\
\texttt{flan-t5-base} ($\approx$250M)  & 0.053 & 0.600 & $+0.547$ &
  83 / 1 / 66 \\
\texttt{flan-t5-large} ($\approx$780M) & 0.113 & 0.633 & $+0.520$ &
  81 / 3 / 66 \\
\bottomrule
\end{tabular}
\end{table}

\begin{table}[h]
\centering
\caption{Exp.~8 --- External distractor count sweep (Dense, chunk\,=\,128,
  $k\!=\!5$, $n\!=\!750$; corpus sizes are approximate).}
\small
\begin{tabular}{rcccc}
\toprule
Distractors & Corpus docs & EM & Token F1 & Recall@5 \\
\midrule
0    & $\sim$2{,}800 & 0.600 & 0.667 & 0.865 \\
500  & $\sim$3{,}300 & 0.599 & 0.667 & 0.865 \\
2000 & $\sim$4{,}800 & 0.600 & 0.668 & 0.865 \\
5000 & $\sim$7{,}800 & 0.600 & 0.667 & 0.865 \\
\bottomrule
\end{tabular}
\end{table}

\section{Statistical Significance Tests}

Table~\ref{tab:sig} reports all nine pairwise contrasts with raw and
Holm-Bonferroni-adjusted $p$-values.
Four contrasts survive correction (bold):
RAG vs.\ No-RAG, both Forced-Miss contrasts, and Oracle vs.\ RAG.

\begin{table}[h]
\centering
\caption{Paired significance results ($n\!=\!750$). McNemar's exact
  test on EM; paired bootstrap on F1 deltas (10{,}000 resamples).
  Holm--Bonferroni correction across all nine contrasts. Discordant
  counts are (A better / B better).}
\label{tab:sig}
\small
\setlength{\tabcolsep}{4pt}
\begin{tabular}{lccccccc}
\toprule
Contrast (A $-$ B) & $\Delta$EM & Disc. &
  McN raw & \textbf{McN Holm} &
  $\Delta$F1 & 95\% CI & \textbf{F1 Holm} \\
\midrule
BM25 $-$ Dense & $-0.003$ & 70/72 & 0.933 & 1.000 &
  $+0.002$ & [$-0.026$, $+0.030$] & 1.000 \\
TF-IDF $-$ Dense & $+0.003$ & 67/65 & 0.931 & 1.000 &
  $+0.005$ & [$-0.021$, $+0.032$] & 1.000 \\
BM25 $-$ TF-IDF & $-0.005$ & 38/42 & 0.738 & 1.000 &
  $-0.003$ & [$-0.025$, $+0.018$] & 1.000 \\
Instructed $-$ Concise & $-0.016$ & 8/20 & 0.036 & 0.179 &
  $-0.012$ & [$-0.025$, $+0.001$] & 0.312 \\
\textbf{RAG $-$ No-RAG} & $\mathbf{+0.551}$ & \textbf{419/6} &
  $\mathbf{<0.001}$ & $\mathbf{<0.001}$ &
  $+0.568$ & [$+0.534$, $+0.603$] & $<0.001$ \\
Dense+Rerank $-$ e5-Dense & $+0.017$ & 47/34 & 0.182 & 0.728 &
  $+0.022$ & [$-0.001$, $+0.044$] & 0.312 \\
\textbf{Oracle $-$ RAG} & $\mathbf{+0.019}$ & \textbf{14/0} &
  $\mathbf{0.0001}$ & $\mathbf{0.0007}$ &
  $+0.016$ & [$+0.007$, $+0.026$] & $0.001$ \\
\textbf{ForcedMiss $-$ No-RAG} & $\mathbf{+0.110}$ & \textbf{104/23} &
  $\mathbf{<0.001}$ & $\mathbf{<0.001}$ &
  $+0.148$ & [$+0.118$, $+0.179$] & $<0.001$ \\
\textbf{ForcedMiss $-$ InstrNoCtx} & $\mathbf{+0.102}$ & \textbf{104/29} &
  $\mathbf{<0.001}$ & $\mathbf{<0.001}$ &
  $+0.145$ & [$+0.114$, $+0.177$] & $<0.001$ \\
\bottomrule
\end{tabular}
\end{table}

\section{Implementation Details}

\paragraph{Chunking.}
Fixed word-window chunking with 32-word overlap.
For a document of $W$ words and chunk size $c$:
chunk $i$ spans words $[i(c-32),\; i(c-32)+c]$.

\paragraph{Generation.}
Flan-T5-base with greedy decoding (\texttt{do\_sample=False},
\texttt{num\_beams=1}, \texttt{max\_new\_tokens=32}).
Long prompts are \emph{middle-truncated}: prefix (question) and suffix
(answer cue) are preserved; only the context body is trimmed.
All generation is cached to
\texttt{data/cache/generation\_cache.json} keyed by MD5 of
\texttt{model\_name + rendered\_prompt}, so different generator sizes
(Exp.~12) never collide.

\paragraph{Dense retrieval.}
Query and passage embeddings are L2-normalised; retrieval is by
dot product (equivalent to cosine similarity).
Passage embeddings are cached per chunk-size to avoid re-encoding across
experiments.
For Exp.~10, each model receives its recommended asymmetric
query/passage instruction prefixes
(\texttt{e5}: \texttt{"query:"}/\texttt{"passage:"}; \texttt{bge}: search
instruction on the query side; \texttt{MiniLM}: no prefix).

\paragraph{Reproducibility.}
All random processes use fixed seeds (corpus sampling seed 42,
bootstrap seed 42).
Running the notebook twice produces identical results.
Expected end-to-end runtime at $n\!=\!750$: approximately 4--8~h on CPU
(significantly less on GPU/MPS with embedding caching).

\end{document}
